% Part: turing-machines
% Chapter: machines-computations
% Section: well-behaved-machines

\documentclass[../../../include/open-logic-section]{subfiles}

\begin{document}

\olfileid{tur}{mac}{dis}
\olsection{শৃঙ্খলাবদ্ধ যন্ত্র}

\begin{explain}
\olref[hal]{sec} অংশে আমরা এমন টুরিং যন্ত্র বিবেচনা করেছি, যার একটিমাত্র
নির্দিষ্ট থামা-দশা~$h$ আছে—এমন যন্ত্র থামলে দশা~$h$-তেই থামবে, তা
নিশ্চিত। এই অর্থে একটিমাত্র থামা-দশাযুক্ত যন্ত্র সাধারণ টুরিং যন্ত্রের
জন্য আমাদের অনুমোদিত আচরণের চেয়ে বেশি ``শৃঙ্খলাবদ্ধ''। টুরিং যন্ত্রের
আচরণের উপর আরও কিছু বিধিনিষেধ আরোপ করা যায়। যেমন, টুরিং যন্ত্রকে ঘর~$0$-
এর ফিতার শেষ-চিহ্নটি কখনও মুছতে অথবা ঘর~$0$ থেকে বাঁ দিকে যাওয়ার চেষ্টা
করতে আমরা নিষেধ করিনি। (আমাদের সংজ্ঞায় এমন ক্ষেত্রে হেডটি ঘর~$0$-তেই
থাকে; অন্য সংজ্ঞায় যন্ত্রটি থেমে যায়।) একইভাবে, টুরিং যন্ত্র থামলে
ঘর~$1$ পড়ার সময়ই থামে—এটি ধরে নিতে পারাও কখনও কাম্য।
\end{explain}

\begin{defn}\ollabel{defn:disciplined}
একটি টুরিং যন্ত্র~$M$ \emph{শৃঙ্খলাবদ্ধ}, যদি এবং কেবল যদি
\begin{enumerate}
    \item তার একটিমাত্র নির্দিষ্ট থামা-দশা~$h$ থাকে,
    \item সেটি থামলে ঘর~$1$ পড়ার সময় থামে,
    \item সেটি ঘর~$0$-এর $\TMendtape$ প্রতীকটি কখনও মোছে না, এবং
    \item সেটি ঘর~$0$ থেকে কখনও বাঁ দিকে যাওয়ার চেষ্টা করে না।
\end{enumerate}
\end{defn}

\begin{explain}
আমরা আগেই আলোচনা করেছি, যেকোনো টুরিং যন্ত্রকে একই আচরণসম্পন্ন কিন্তু
নির্দিষ্ট থামা-দশাযুক্ত যন্ত্রে বদলানো যায়। তার জন্য শুধু একটি নতুন
দশা~$h$ যোগ করতে হয়, এবং আদি~$\delta$ যে প্রতিটি জোড়া
$\tuple{q,\sigma}$-তে অসংজ্ঞায়িত, তার জন্য একটি নির্দেশ
$\delta(q, \sigma) = \tuple{h, \sigma, N}$ যোগ করতে হয়। প্রমাণটি
ক্লান্তিকর হলেও সত্য যে, যেকোনো টুরিং যন্ত্র~$M$-কে এমন একটি
শৃঙ্খলাবদ্ধ টুরিং যন্ত্র~$M'$-তে বদলানো যায়, যা একই নিবেশে থামে এবং একই
নির্গম উৎপন্ন করে। যেমন, টুরিং যন্ত্রটি থামলে যদি ঘর~$1$-এ না থাকে, তবে
হেডটি ফিতার শেষ-চিহ্ন না পাওয়া পর্যন্ত বাঁ দিকে সরবে, তারপর এক ঘর ডান
দিকে গিয়ে থামবে—এমন কিছু নির্দেশ যোগ করা যায়। অন্য শর্তগুলি কীভাবে
সামলানো যায়, তা তোমাদের ভেবে দেখতে দিলাম।
\end{explain}

\begin{ex}
    \olref{fig:adder-disc}-এ আমরা \olref[una]{ex:adder}-এর যোগ-যন্ত্রটিকে
    একটি শৃঙ্খলাবদ্ধ যন্ত্রে বদলাই।
    \begin{figure}\[
\begin{tikzpicture}[->,>=stealth',shorten >=1pt,auto,node distance=2.8cm,
                    semithick]
  \tikzstyle{every state}=[fill=none,draw=black,text=black]

  \node[initial,state]         (A)              {$q_0$};
  \node[state]         (B) [right of=A] {$q_1$};
  \node[state]         (C) [below right of=B] {$q_2$};
  \node[state]         (D) [below left of=C] {$q_3$};
  \node[state]         (H) [left of=D] {$h$};

  \path (A) edge node {\TMtrans{\TMblank}{\TMstroke}{\TMstay}} (B)
           edge [loop below] node {\TMtrans{\TMstroke}{\TMstroke}{\TMright}} (A)
        (B) edge [loop below] node {\TMtrans{\TMstroke}{\TMstroke}{\TMright}} (B)
            edge node {\TMtrans{\TMblank}{\TMblank}{\TMleft}} (C)
        (C) edge node {\TMtrans{\TMstroke}{\TMblank}{\TMleft}} (D)
        (D) edge [loop above] node {\TMtrans{\TMstroke}{\TMstroke}{\TMleft}} (D)
        edge node {\TMtrans{\TMendtape}{\TMendtape}{\TMright}} (H);
\end{tikzpicture}
\]\caption{একটি শৃঙ্খলাবদ্ধ যোগ-যন্ত্র}
\ollabel{fig:adder-disc}
\end{figure}
% BN-SRC-170 TARGET-CORRECTION-BEGIN
\par\smallskip\noindent\textnormal{\small\textbf{উৎস-সংশোধন (BN-SRC-170):} হিমায়িত ইংরেজি চিত্রে $q_0$-তে $\TMstroke$ পড়ার তীরটি ভুল করে $q_1$-এ যায়। বর্ণিত যোগ-পদ্ধতি অনুযায়ী বাংলা চিত্রে তীরটি $q_0$-র স্ব-আবর্তন করা হয়েছে; তারপরেই যন্ত্রটি প্রথম খণ্ড পেরিয়ে তার ফাঁকা বিভাজকে $q_1$-এ যায়।} % BN-SRC-170 TARGET-CORRECTION-END
\end{ex}

\begin{prop}\ollabel{prop:disciplined} প্রতিটি টুরিং যন্ত্র~$M$-এর জন্য
এমন একটি শৃঙ্খলাবদ্ধ টুরিং যন্ত্র~$M'$ আছে, যাতে $M$ নির্গম~$O$ নিয়ে
থামলে সেটিও নির্গম~$O$ নিয়ে থামে, আর $M$ না থামলে সেটিও থামে না।
বিশেষত, কোনো টুরিং যন্ত্রে গণনসাধ্য প্রতিটি অপেক্ষক
$f\colon\Nat^n \to \Nat$ একটি শৃঙ্খলাবদ্ধ টুরিং যন্ত্রেও গণনসাধ্য।
\end{prop}

\begin{prob}\label{tur:mac:dis:prob:disc-succ}
$f(x) = x+1$ গণনা করে এমন একটি শৃঙ্খলাবদ্ধ যন্ত্র দাও।
\end{prob}


\begin{prob}\label{tur:mac:dis:prob:copier}
এমন একটি শৃঙ্খলাবদ্ধ যন্ত্র খুঁজে বার করো, যা নিবেশ
$\TMstroke^n$-এ শুরু করে নির্গম
$\TMstroke^n \concat \TMblank \concat \TMstroke^n$ উৎপন্ন করে।
\end{prob}

\end{document}
